% AY2014 制御工学 第1問（転記）
% 出典: 03_制御工学/original/03_制御工学/2014/2014se.pdf
% 要約: 二次振動系（質量 m, 粘性係数 d, バネ定数 k, 力 f(t), 変位 x(t)）。
%       伝達関数・固有振動数・減衰係数、パラメータ決定、ステップ応答、正弦入力応答。

\section*{第1問}

右図の二次振動系に関する次の問い (1)〜(4) に答えよ. なお, 質量を $m$, 粘性係数を $d$,
バネ定数を $k$ とし, $f(t)$ は系に加える力, $x(t)$ は質量中心の位置を表す.

\begin{center}
\begin{tikzpicture}[>=Latex]
  % 壁
  \draw[thick] (0,-0.2) -- (0,2.2);
  \foreach \y in {-0.1,0.1,...,2.05} \draw[thick] (0,\y) -- (-0.25,\y-0.25);
  % バネ k（上）
  \draw[thick,decorate,decoration={zigzag,segment length=8,amplitude=5}] (0,1.6) -- (3,1.6);
  \node at (1.5,2.05) {$k$};
  % ダンパ d（下）
  \draw[thick] (0,0.45) -- (1.15,0.45);
  \draw[thick] (1.15,0.2) rectangle (1.85,0.7);
  \draw[thick] (1.55,0.45) -- (3,0.45);
  \node at (1.5,-0.05) {$d$};
  % 質量 m
  \draw[thick,fill=gray!12] (3,0.1) rectangle (3.9,2.0);
  \node at (3.45,1.05) {$m$};
  % 力 f(t)
  \draw[->,thick] (3.9,1.05) -- (5.0,1.05);
  \node at (5.25,1.05) {$f(t)$};
  % 変位 x(t)
  \draw[->,thick] (3.45,2.35) -- (4.3,2.35);
  \node at (3.9,2.65) {$x(t)$};
\end{tikzpicture}
\end{center}

\begin{enumerate}[label=(\arabic*)]
  \item 入力を $f(t)$, 出力を $x(t)$ として, システムの伝達関数を求めたうえで,
        固有振動数ならびに減衰係数を求めよ.

  \item 物理パラメータを $m=1$, $d=1$, $k=\beta$, 初期値を $x(0)=5$, $\dot{x}(0)=1$ とし,
        入力 $f(t)$ を単位ステップ関数とする. このとき, $x(t)$ の定常値が $5$ になるように
        $\beta$ を定めよ.

  \item 物理パラメータを $m=1$, $d=3$, $k=2$, 初期値を $x(0)=5$, $\dot{x}(0)=-1$ とし,
        入力を単位ステップ関数とするときの $x(t)$ を求めよ. 加えて, $x(t)$ の概形を示せ.
        なお, 横軸を時間, 縦軸を $x(t)$ とすること.

  \item 物理パラメータを $m=1$, $d=1$, $k=1$, 初期値を $x(0)=0$, $\dot{x}(0)=0$ とし,
        入力を $f(t)=\sin(\omega t)$ とする（ただし, $t<0$ のとき $f(t)=0$）.
        このときの $x(t)$ を求めよ.
\end{enumerate}
